\documentclass{article}

\usepackage[utf8]{inputenc}     % pdfLaTeX
\usepackage[T1]{fontenc}
\usepackage[magyar]{babel}

\title{A Középkor Irodalma}
\author{Haydn Lorenzo}
\date{Legutóbb frissítve: 2026. Szeptember 18.}

\begin{document}

\maketitle
\tableofcontents
\newpage

\section{Rímek}
\begin{itemize}
    \item Elhelyezkedés szerint:
    \begin{itemize}
        \item végrím
        \item füzérrím "[...] istenek felett. Ég kételkedett [...]
        \item sorrím "Párja várja, hogy álma felrázza"
        \item belső rím "Monyók álma a hétfői matek s, csakis azt várja? hogy elmúljon"
        \item betűrím (alliteráció)
    \end{itemize}
    \item Rímképletek szerint
    \begin{itemize}
        \item bokorrím (aaaa)
        \item párosrím (aabb)
        \item ölelkező rím (abba)
        \item keresztrím (abab)
        \item tiszta rím vagy önrím (mgh. azonos)
        \item asszonánc $\rightarrow$ mgh. és msh. azonos pl.: zománc, asszonánc
    \end{itemize}
\end{itemize}



\end{document}